

Ce manuel est composé de l'ensemble des activités, cours, exercices pour les classes de 5\up{e} du collège Simone Veil de Montpellier que j'ai à ma charge durant l'année 2022-2023. \\ [1mm]
Il a été écrit en \LaTeX{} avec la classe \href{https://www.ctan.org/pkg/sesamanuel}{\bf Sesamanuel} distribuée librement par l'association Sesamath, ainsi que le package \href{https://www.ctan.org/pkg/ProfCollege}{\bf ProfCollege}. Une partie des exercices est issue du \og Manuel Sesamath, 5\up{e} \fg{} et du \og Cahier Sesamath, 5\up{e} \fg, éditions génération 5 et Magnard-Sesamath.

Si vous y voyez des erreurs ou des coquilles, même minimes, vous pouvez me les signaler à cette adresse : \href{mailto:nathalie.daval@ac-montpellier.fr}{nathalie.daval@ac-montpellier.fr}. Je remercie à ce propos {\it Jean-Félix Navarro} qui a effectué une relecture attentive du livret précédent, {\it Sébastien Lozano} pour son aide précieuse en \LaTeX{} et {\it Christophe Poulain} pour son package ProfCollege et les améliorations au fil de l'eau.  \\ [10mm]

La progression est dite spiralée, c'est-à-dire que chaque \og chapitre \fg{} est décomposé en plusieurs séquences conçues pour durer une semaine en moyenne, ce qui permet de revoir les notions plusieurs fois dans l'année. La page suivante propose une programmation possible sur les cinq périodes (P1, P2, P3, P4 et P5) de l'année 2022-2023.

Chaque séquence du présent manuel est composée de la manière suivante : \\
\begin{itemize}
   \item \colorbox{J4}{\textcolor{B2}{\sffamily\bfseries Ce que sait faire l'élève en 5\up{e}}} : les attendus de fin d'année de cinquième en mathématiques, tels qu'ils sont définis par le \href{https://eduscol.education.fr/document/14044/download}{document Eduscol}. Ils serviront de base pour les évaluation des compétences. \medskip
   \item \textcolor{B2}{\sffamily\bfseries Connaissances et compétences associées} : les connaissances et compétences associées au cycle 4 définies par le \href{https://eduscol.education.fr/document/621/download}{programme en vigueur à compter de la rentrée de l'année scolaire 2020}. \medskip
   \item \textcolor{B1}{\sffamily\bfseries Débat} : un petit texte culturel illustré permettant d'échanger sur un thème en rapport au chapitre. Un morceau d'histoire, de l'étymologie, du vocabulaire, une curiosité mathématique\dots{} le tout agrémenté d'une courte vidéo de vulgarisation scientifique. \medskip
   \item \colorbox{A2}{\textcolor{white}{\sffamily\bfseries Activité d'approche}} : une activité à faire en classe, permettant de découvrir une notion du chapitre et de construire la trace écrite. \medskip
   \item \colorbox{A1}{\textcolor{white}{\sffamily\bfseries Trace écrite}} : l'essentiel du cours à connaître. \medskip
   \item \colorbox{B1}{\textcolor{white}{\sffamily\bfseries Entraînement}} : une fiche d'exercices suivant les connaissances et compétences à acquérir. \medskip
   \item \colorbox{B2}{\textcolor{white}{\sffamily\bfseries Récréation}} : une activité ludique liée à la séquence, une énigme ou un problème à résoudre.
   \item
\end{itemize}

\bigskip

En fin de manuel se trouvent les {\sffamily PLANS DE TRAVAIL} donnés aux élèves, pour chaque notion. Chacun de ces plans est constitué de : \\
   $\bullet$ la bulle \og Je connais mon cours \fg{} qui récapitule les compétences à travailler ; \\
   $\bullet$ la bulle \og Aide en vidéo \fg{} composée de deux capsules issues sur site d'{\it Yvan Monka} : \href{https://www.maths-et-tiques.fr}{m@ths et tiques} ; \\
   $\bullet$ la bulle \og Questions flash \fg{} qui permet de noter le résultat aux questions flash de la séquence ; \\
   $\bullet$ des bulles de compétences, celles-ci indiquent les exercices associés à chaque compétence, ainsi que son niveau de difficulté ($\star$, $\star\star$ ou $\star\star\star$).
   
\vfill

\pagebreak

\newcommand\gm[2]{\rput(2,#1){\begin{minipage}{6cm}\centering\textcolor{Green}{#2}\end{minipage}}}
\newcommand\nc[2]{\rput(6,#1){\begin{minipage}{7cm} \centering\textcolor{Red}{#2}\end{minipage}}}
\newcommand\eg[2]{\rput(10,#1){\begin{minipage}{6cm}\centering\textcolor{DodgerBlue}{#2}\end{minipage}}}
\newcommand\ogd[2]{\rput(14,#1){\begin{minipage}{6cm}\centering\textcolor{violet}{#2}\end{minipage}}}

{\psset{xunit=1,yunit=0.7}


{\small
\begin{pspicture}(-0.5,0)(17,33)
   \rput(0.5,35){\includegraphics[width=1cm]{Images/gm}}
   \gm{35}{\bf Grandeurs\\\& mesures}
   \rput(4.5,35){\includegraphics[width=1cm]{Images/nc}}
   \nc{35}{\bf Nombres\\\& calculs}
   \rput(8.5,35){\includegraphics[width=1cm]{Images/eg}}
   \eg{35}{\bf Espace \&\\géométrie}
   \rput(12.3,35){\includegraphics[width=1cm]{Images/ogd}}
   \ogd{35}{\bf Organisation\\\& gestion\\de données}
   \multido{\n=-1+1}{35}{\psline[linestyle=dotted,linecolor=gray](0,\n)(16,\n)}
   \psline(-1,34)(17,34)
   \rput(-0.5,30){$P_1$}
   \psline(-1,26)(17,26)
   \rput(-0.5,23){$P_2$}
   \psline(-1,20)(17,20)
   \rput(-0.5,16.5){$P_3$}
   \psline(-1,13)(17,13)
   \rput(-0.5,9.5){$P_4$}
   \psline(-1,6)(17,6)
   \rput(-0.5,2.5){$P_5$}
   %%%%%%%%% P1
   \rput(8,33.5){\gray Semaine de rentrée}
   \nc{32.5}{Nombres et calculs 1\\1. Enchaînement d'opérations}
   \eg{31.5}{Géométrie plane, démonstrations 1\\2. Angles particuliers}
   \rput(8,30.5){\textcolor{orange}{3. En route vers la programmation ({\it Introduction... puis fil rouge tout au long de l'année})}}
   \nc{29.5}{Nombres et calculs 2\\4. Nombres relatifs}
   \eg{28.5}{Représenter l'espace 1\\5. Repérage dans le plan}
   \ogd{27.5}{Statistiques 1\\6. Interpréter, représenter des données}
   \rput(8,26.5){\gray Semaine de rattrapage 1}
   %%%%%%%%% P2
   \gm{25.5}{Grandeurs mesurables 1\\7. Horaires et durées}
   \nc{24.5}{Calcul littéral 1\\8. Expressions littérales}
   \eg{23.5}{Géométrie plane, démonstrations 2\\9. Somme des angles d'un triangle}
   \ogd{22.5}{Probabilités 1\\10. Notions de probabilités}
   \nc{21.5}{Arithmétique 1\\11. Multiples et diviseurs}
   \rput(8,20.5){\gray Semaine de rattrapage 2}
   %%%%%%%%% P3
   \eg{19.5}{Géométrie plane, démonstrations 3\\12. La symétrie centrale}
   \gm{18.5}{Grandeurs mesurables 2\\13. Calcul d'aires}
   \nc{17.5}{Nombres et calculs 3\\14. Comparaison et égalité de fractions}
   \eg{16.5}{Géométrie plane, démonstrations 4\\15. L'inégalité triangulaire}
   \ogd{15.5}{Proportionnalité 1\\16. Proportionnalité}
   \nc{14.5}{Calcul littéral 2\\17. Distributivité simple et égalités}
   \rput(8,13.5){\gray Semaine de rattrapage 3}
   %%%%%%%%% P4  
   \eg{12.5}{Représenter l'espace 2\\18. Reconnaître des solides}
   \gm{11.5}{Grandeurs mesurables 3\\19. Volume du prisme et du cylindre} 
   \nc{10.5}{Nombres et calculs 4\\20. Somme et différence de nombres relatifs}
   \eg{9.5}{Géométrie plane, démonstrations 5\\21. Le parallélogramme}
   \ogd{8.5}{Proportionnalité 2\\22. Le ratio}
   \nc{7.5}{Arithmétique 2\\23. Nombres premiers}
   \rput(8,6.5){\gray Semaine de rattrapage 4}
   %%%%%%%%% P5
   \eg{5}{Représenter l'espace 3\\24. Représenter des solides}
   \gm{3.5}{Grandeurs mesurables 4\\25. L'aire du parallélogramme}
   \nc{2.5}{Nombres et calculs 5\\26. Somme et différence de fractions}
   \ogd{1.5}{Statistiques 2\\27. Fréquence et moyenne}
   \eg{0.5}{Géométrie plane, démonstrations 6\\28. Les droites du triangle}
   \gm{-0.5}{Effet des transformations 1\\29. Propriétés des symétries}
\end{pspicture}}}


\pagebreak

\textsf{\huge SOMMAIRE}

%bandeau nombres et calculs
\begin{pspicture}(0,0)(\linewidth,1.5)
   \psframe*[linecolor=Red](0,0)(\linewidth,1)
   \rput(8.4,0.5){\textcolor{white}{\Large\textsf{NOMBRES ET CALCULS}}}
\end{pspicture}

\begin{multicols}{2}
   S01 Enchaînement d'opérations \pointilles \pageref{S01} \\
   S04 Nombres relatifs \pointilles \pageref{S04} \\
   S08 Expressions littérales \pointilles \pageref{S08} \\
   S11 Multiples et diviseurs \pointilles \pageref{S11} \\
   S14 Comparaison et égalité de fractions \pointilles \pageref{S14} \\
   S17 Distributivité simple et égalités \pointilles \pageref{S17} \\
   S20 Somme et différence de nombres relatifs \pointilles \pageref{S20} \\
   S23 Nombres premiers \pointilles \pageref{S23} \\
   S26 Somme et différence de fractions \pointilles \pageref{S26} \\
\end{multicols}
   
%bandeau géométrie
\begin{pspicture}(0,0)(\linewidth,1.5)
   \psframe*[linecolor=DodgerBlue](0,0)(\linewidth,1)
   \rput(8.4,0.5){\textcolor{white}{\Large\textsf{GÉOMÉTRIE}}}\end{pspicture}
    
\begin{multicols}{2}
   S02 Angles particuliers \pointilles \pageref{S02} \\
   S05 Repérage dans le plan \pointilles \pageref{S05} \\
   S09 Somme des angles d'un triangle \pointilles \pageref{S09} \\
   S12 La symétrie centrale \pointilles \pageref{S12} \\
   S15 L'inégalité triangulaire \pointilles \pageref{S15} \\
   S18 Reconnaître des solides \pointilles \pageref{S18} \\      
   S21 Le parallélogramme \pointilles \pageref{S21} \\
   S24 Représenter des solides \pointilles \pageref{S24} \\
   S28 Les droites du triangle \pointilles \pageref{S27}   
\end{multicols}
   
%bandeau organisation et gestion de données
\begin{pspicture}(0,0)(\linewidth,1.5)
   \psframe*[linecolor=violet](0,0)(\linewidth,1)
   \rput(8.4,0.5){\textcolor{white}{\Large\textsf{ORGANISATION ET GESTION DE DONNÉES}}}
\end{pspicture} 

\begin{multicols}{2}
   S06 Interpréter et représenter des données \pointilles \pageref{S06} \\
   S10 Notions de probabilités \pointilles \pageref{S10} \\
   S16 Proportionnalité \pointilles \pageref{S16} \\
   S22 Le ratio \pointilles \pageref{S22} \\
   S27 Fréquence et moyenne \pointilles \pageref{S28}
\end{multicols}

%bandeau grandeurs et mesures
\begin{pspicture}(0,0)(\linewidth,1.5)
   \psframe*[linecolor=Green](0,0)(\linewidth,1)
   \rput(8.4,0.5){\textcolor{white}{\Large\textsf{GRANDEURS ET MESURES}}}\end{pspicture}
    
\begin{multicols}{2}
   S07 Horaires et durées \pointilles \pageref{S07} \\
   S13 Calcul d'aires \pointilles \pageref{S13} \\
   S19 Volume du prisme et du cylindre \pointilles \pageref{S19} \\
   S25 L'aire du parallélogramme \pointilles \pageref{S25} \\
   S29 Propriétés des symétries \pointilles \pageref{S29}
\end{multicols}
   
%bandeau programmation
\begin{pspicture}(0,0)(\linewidth,1.5)
   \psframe*[linecolor=orange](0,0)(\linewidth,1)
   \rput(8.4,0.5){\textcolor{white}{\Large\textsf{ALGORITHMIQUE ET PROGRAMMATION}}}
\end{pspicture}
    
\begin{multicols}{2}
   S03 En route vers la programmation \pointilles \pageref{S03} \\
\end{multicols}

Plans de travail \pointilles \pageref{PDT}

Solution des exercices et de la partie récréation \pointilles \pageref{solutions}
